\documentclass[11pt,a4paper]{article}
\usepackage[margin=27mm]{geometry}
\usepackage{amsmath,amsthm,amssymb,fontspec,unicode-math,hyperref}
\usepackage{longtable,booktabs,array,calc,fvextra}
\setmainfont{STIXGeneral}
\setmathfont{STIX Two Math}
\setmonofont{Menlo}
\hypersetup{hidelinks}
\newtheorem{theorem}{Theorem}
\newtheorem{proposition}[theorem]{Proposition}
\newtheorem{lemma}[theorem]{Lemma}
\newtheorem{corollary}[theorem]{Corollary}
\theoremstyle{remark}\newtheorem{remark}[theorem]{Remark}
\newcommand{\gen}[1]{\langle #1\rangle}
\newcommand{\F}{\mathbb F}
\newcommand{\Z}{\mathbb Z}
\newcommand{\Aut}{\operatorname{Aut}}
\newcommand{\Sym}{\operatorname{Sym}}
\setlength{\emergencystretch}{2em}
\providecommand{\tightlist}{\setlength{\itemsep}{0pt}\setlength{\parskip}{0pt}}

\title{Petrials and faithful chamber actions}
\author{Human Compression Project}
\date{}
\begin{document}
\maketitle
\begin{abstract}
The geometric Petrial of a thin regular rank-three geometry retains
vertices and edges and identifies equal Petrie sets. We compute the
kernel of its natural automorphism action and prove that a thin
Petrial forces a commuting pair of chamber generators. Under that
condition the induced action is faithful and regular, and taking the
Petrial twice recovers the original geometry. The proof covers
infinite geometries and infinite residues. Two explicit examples
separate these conclusions from unconditional embedding and
length assertions in the original article.
\end{abstract}
\section{The source and common prerequisites}
The target is Bonheure--Buekenhout--Leemans [BBL01, Theorem 1], together
with the specified definition and Lemma 3.2 conclusions. We use the
complete authored text rendering of its seven printed pages 19--25;
the original PDF was unavailable. This edition supplies the statement
boundary, not a PDF length for comparison.

The common paper [U] proves connected chamber galleries, their marked
stabilizers and the exact kernel when retained objects determine the
new ones. These statements require neither a finite group nor finite
residues. Here the new faces are actual pairs of vertex and edge sets,
not independently named walks. This distinction is essential to the
kernel calculation.
\section{Geometric Petrials and the exact faithful
image}\label{geometric-petrials-and-the-exact-faithful-image}

The operation in this paper is that of Bonheure--Buekenhout--Leemans
(2001), with equal geometric Petrie sets identified. Its principal
theorem and valid numbered structural conclusions are the targets. The
unconditional embedding assertion is corrected below. No source
conclusion is used as a hypothesis, and neither a finite group nor
finite residues are assumed.

Let \(\Gamma\) be a thin regular residually connected rank-three
geometry. Label its types vertex, edge and face. At a base chamber
\(C=(v,e,f)\), let a change the face, b the edge and c the vertex. Lemma
C.Gallery gives

\[
 G=\operatorname{Aut}(\Gamma)=\langle a,b,c\rangle,
 \quad G_v=\langle a,b\rangle,\quad G_e=\langle a,c\rangle,
 \quad G_f=\langle b,c\rangle,\quad G_v\cap G_e=\langle a\rangle.
 \tag{P1}
\]

The other two pair intersections are respectively \(\langle b\rangle\)
and \(\langle c\rangle\), and the chamber stabilizer is trivial. Any two
of a,b,c generate a dihedral group, possibly infinite: a word in two
involutions alternates after cancellation and is a power of their
product or such a power times one involution. For distinct nonidentity
involutions x,y, the intersection
\(\langle xy\rangle \cap \langle x\rangle\) is trivial. Otherwise x is a
power of xy and commutes with it, forcing \(xy=yx\) and \(|xy|=2\); then
\(x=xy\) would force \(y=1\). This proves the rotation/reflection fact
used below without a finiteness assumption.

\subsection{The operation and its faithful
image}\label{the-operation-and-its-faithful-image}

Put \(w=cab\). Starting at a chamber, change vertex, face, then edge; at
the base chamber the result is wC. At gC, successive chamber flips act
by right multiplication on g, so iteration gives \(gw^iC\). The Petrie
face at C is the actual ordered pair

\[
 P(C)=\left(\{w^iv:i\in\mathbb Z\},\{w^ie:i\in\mathbb Z\}\right).
 \tag{P2}
\]

Equal pairs are one face; the vertex--edge incidence induced on these
sets is understood. Define \(\Pi \Gamma\) by retaining old vertices,
edges and their incidence, and adjoining these distinct faces, with
incidence to a face given by membership. This is equivariant:
\(P(gC)=gP(C)\). In particular \(wP(C)=P(C)\), and

\[
 wv=cv,\qquad w^{-1}e=be.
 \tag{P3}
\]

Thus a Petrie face contains v and cv at e, and e and be at v. Each pair
is distinct because the corresponding old panel change is nontrivial.

\textbf{Lemma C.PetrialAction (connectivity and exact action).}
\(\Pi \Gamma\) is a residually connected geometry. Its natural
automorphism action \(\rho\) has kernel

\[
 \ker\rho=\operatorname{core}_G\langle a\rangle
 =\begin{cases}\langle a\rangle,&a\in Z(G),\\1,&a\notin Z(G).\end{cases}
 \tag{P4}
\]

\emph{Proof.} Every incident old vertex--edge pair extends through an
old chamber to its Petrie face. For an object belonging to a Petrie
face, a step of the defining walk supplies an incident object of the
other retained type. These facts extend every flag to a chamber.

Write \(v_i=w^iv\) and \(e_i=w^ie\). Equation (P3) says that \(e_i\)
joins \(v_i\) and \(v_{i+1}\). The walk therefore connects every object
in its face residue; additional induced incidences preserve
connectivity. In the old residue at v, replace every two-step path
\(e_{1}–f–e_{2}\) with distinct edges by \(e_{1}–P(v,e_{1},f)–e_{2}\).
Thinness and (P3) justify the replacement. It connects all retained
edges in the new vertex residue, and every new incident face meets one
of them. The analogous replacement \(v_{1}–f–v_{2}\) proves connectivity
of each edge residue. Finally an old path through a face can be replaced
by a path in its connected vertex--edge residue. Thus the old
vertex--edge truncation, and consequently the new geometry with every
face attached, is connected.

For the action apply UC.Action with X the disjoint union of old vertices
and edges. They are retained in Y, and every new face is uniquely
encoded by its two actual subsets of X. UC.Action's core formula and
(P1) give the core in (P4). A normal subgroup of \(\langle a\rangle\) is
either trivial or all of \(\langle a\rangle\). The latter is normal in G
exactly when its unique nonidentity element a is central. This proves
the exact formula and constructs the faithful image \(G/ker\rho\).
\(\square\)

This is the same retained-object theorem used for graph halving in the common paper,
but the smaller retained object set can have a nontrivial pointwise
kernel. It asserts neither chamber transitivity of \(\Pi \Gamma\) nor
injectivity of \(\rho\) without further proof.

For an ordinary map, the three successive changes move across an edge to
its other endpoint, change the incident face, and continue along that
face. These are exactly the alternating corner choices in the
conventional Petrie walk. This proves the walk-definition assertion of
source Lemma 3.1; it does not forbid repetitions or replace geometric
faces by labelled walks.

\subsection{Thinness and recovery}\label{thinness-and-recovery}

\textbf{Theorem C.Petrial (Petrial obstruction and recovery).} If
\(\Pi \Gamma\) is thin, then \(|ab|=2\) or \(|ac|=2\). The action
\(\rho\) is faithful and regular and is the full automorphism group of
\(\Pi \Gamma\). Applying \(\Pi\) again recovers \(\Gamma\) by an
isomorphism fixing all retained vertices and edges.

\emph{Proof.} Set \(P=P(C)\). Suppose first that \(bP\ne P\). Both P and
bP contain e and be at v, so the vertex residue of \(\Pi \Gamma\)
contains the four-cycle e--P--be--bP--e. A finite cycle in a connected
graph of degree two is the entire graph: its vertices already have their
two neighbours, so no path can leave it. Lemma C.PetrialAction therefore
says that precisely two old edges are incident with v. In the old thin
connected vertex residue both edges have degree two, and each incident
face meets both. It is also a four-cycle. Its regular dihedral group
\(\langle a,b\rangle\) has four chambers, so \(|ab|=2\).

Suppose instead that \(bP=P\). Since w fixes P, \(ca=wb\) fixes P and e.
By the general freeness statement of UC.Geometry, the automorphism group
of the thin \(\Pi \Gamma\) acts freely on chambers. Its edge--face panel
stabilizer embeds in the permutation group on that panel's two
completions, so \(\rho (ca)^2=1\). Equation (P4) puts \((ca)^2\) in
\(\langle a\rangle\). The rotation/reflection intersection in
\(\langle a,c\rangle\) is trivial, hence \((ca)^2=1\). Since \(c\ne a\),
this gives \(|ac|=2\). This proves the obstruction without presuming
that b always fixes P or that the new action is regular.

Interchanging vertex and edge interchanges b,c and changes cab to
\(bac=w^{-1}\), preserving the same Petrie sets after that type
interchange. We may consequently assume \([a,c]=1\). Then
\(bwb=w^{-1}\), \(bv=v\) and \(be=w^{-1}e\), so \(bP=P\). The subgroup
\(\langle b,w\rangle =\langle b,ca\rangle\) is transitive on the
incident vertex--edge pairs of P: its elements give \((v_i,e_i)\) and
\((v_i,e_{i-1})\) for every i. These are distinct pairs, and thinness
makes them all pairs at \(v_i\). Every object of P occurs in the walk.
The original chamber-transitive G is transitive on the faces P(gC), so
its induced action on \(\Pi \Gamma\) is chamber-transitive. UC.Action
identifies its image with the full regular automorphism group. The
preimage of the new panel stabilizing (v,e) is exactly
\(G_v\cap G_e=\langle a\rangle\). UC.Action's panel-index formula forces
\(ker\rho =1\).

At (v,e,P), the new panel involutions are therefore a,b,ca: the last
fixes (e,P) and sends v to cv; b fixes (v,P) and changes e; a exchanges
the two faces at (v,e). The new generators generate the same group since
\(c=(ca)a\). Their geometric marking is established by this action, not
by assuming a homomorphism from the old marking. Their second Petrie
step is \((ca)ab=cb\). The cb-orbits of v and e exhaust exactly the old
face's vertex--edge shadow: use the alternating normal form in its
residue group \(\langle b,c\rangle\) from (P1).

Those old face shadows are distinct. Otherwise let \(f\ne f'\) have the
same shadow and choose (v,e) in it, moving the base chamber if needed.
Thinness gives \(f'=af\). Both \(f,f'\) meet e and be at v. Their
four-cycle exhausts the old vertex residue, as above, forcing
\([a,b]=1\). Together with \([a,c]=1\) this makes a central in G, and
(P4) contradicts \(ker\rho =1\). Thus old faces correspond bijectively
to the second Petrie faces, with precisely their original incidences.
This proves the required isomorphism. The vertex--edge-swapped case
gives the same conclusion. \(\square\)

The commuting pair is the source's linear-diagram conclusion. Theorem
C.Petrial proves its Theorem 1 and the thin-only formulation in its
abstract, together with Lemma 3.2(iii),(iv). Lemma C.PetrialAction gives
Lemma 3.2(i). No finite-residue assumption occurs in the walk,
connected-cycle, dihedral or gallery arguments.

\subsection{Corrected assertions}\label{corrected-assertions}

The unconditional Lemma 3.2(ii) must be read as the natural homomorphism
with kernel (P4), not as an embedding. Indeed take the complete
tripartite geometry with two objects of each type. Its full
type-preserving group is \(C_{2}^{3}\), regularly permuting eight
chambers. Let a,b,c independently change the face, edge and vertex
coordinates. The element cab changes all three, so every Petrie face
consists of both vertices and both edges. Deduplication leaves one face.
Thus \(\Pi \Gamma\) is the complete tripartite geometry with type sizes
(2,2,1), its full group is \(C_{2}^{2}\), and the kernel is
\(\langle a\rangle\). An order-eight group cannot embed in this
order-four automorphism group. This also gives a commuting-pair diagram
with nonthin Petrial, so the principal implication has no converse in
this diagram convention.

The unnumbered geometric length/order equality also cannot be
unconditional. Take \(C_{2}\times S_{3}\), with a generating its central
\(C_{2}\) and b,c distinct transpositions generating \(S_{3}\). The
coset geometry with stabilizers in (P1) is a six-cycle on three vertices
and three edges, joined to two faces. It is thin, residually connected
and regular: its twelve chambers form the regular \(C_{2}\times S_{3}\)
action. On retained objects, cab acts as cb, a rotation of order three,
while cab itself has order six. Its geometric Petrie face has three
vertices and three edges, not length six. This assertion is unnecessary
for every proof above.

Both counterexamples are proved directly. A separately supplied optional
program checks seven finite examples using literal chamber changes and
deduplicated vertex--edge sets; it is not a premise for the theorem or
an exhaustive classification. The uniform proof above covers infinite
groups and residues. No source atlas or external map classification is
needed.

\section*{Attribution and formal scope}
The Petrial operation and its principal thinness obstruction belong to
[BBL01]. The proof here reuses the project's accepted complete argument,
including the exact action kernel and its written counterexamples.
The common retained-object theorem replaces a repeated action argument;
the obstruction and recovery steps remain specific to Petrie walks.
No historical priority for the corrections is asserted.

The principal theorem has no complete source-to-Lean instantiation.
Generic action or subgroup components must not be described as a
formal certificate of its infinite geometric proof. The finite
diagnostic examples in the earlier work are optional checks, not
premises for any conclusion here.

\section*{References}
{[}U{]} \emph{Changes of chamber generators and binary translation
groups}, the common paper of Topic C.

{[}BBL01{]} D. Bonheure, F. Buekenhout and D. Leemans,
\emph{On the Petrials of thin rank 3 geometries}, Journal of Geometry
71 (2001), 19--25. DOI: 10.1007/s00022-001-8548-2.
\end{document}
